% !TeX program = xelatex
% !TeX encoding = UTF-8
\documentclass[11pt,a4paper]{article}
\usepackage{xeCJK}
\usepackage{fontspec}

% 한글 글꼴 설정 (시스템에 있는 글꼴)
\setCJKmainfont{Adobe Myungjo Std}
\setCJKsansfont{Malgun Gothic}
\setCJKmonofont{Malgun Gothic}

% 라틴 문자 글꼴 설정
\setmainfont{Times New Roman}
\setmonofont{Courier New}

\usepackage{amsmath,amssymb,amsthm,mathtools}
\usepackage{geometry}
\usepackage{booktabs}
\usepackage{hyperref}
\usepackage{url}
\usepackage{tabularx}
\usepackage{array}
\usepackage{makecell}
\usepackage{ragged2e}
\usepackage{bm}
\usepackage{slashed}
\usepackage{tikz}
\usetikzlibrary{arrows.meta, positioning, calc, shapes.geometric}

\usepackage{graphicx}
\usepackage{caption}
\usepackage{subcaption}
\usepackage{float}
\usepackage{titlesec}

\geometry{left=1.5cm,right=1.5cm,top=1.5cm,bottom=2.0cm}

% ============================================================
% 정리 환경 (한국어)
% ============================================================
\newtheorem{theorem}{정리}[section]
\newtheorem{axiom}[theorem]{공리}
\newtheorem{lemma}[theorem]{보조정리}
\newtheorem{corollary}[theorem]{따름정리}
\newtheorem{proposition}[theorem]{명제}
\newtheorem{definition}[theorem]{정의}
\theoremstyle{remark}
\newtheorem{remark}[theorem]{비고}
\renewenvironment{proof}{\paragraph{증명:}}{\hfill$\square$}

% ============================================================
% YUCT 매크로
% ============================================================
\newcommand{\Keff}{K_{\mathrm{eff}}}
\newcommand{\Keffcrit}{K_{\mathrm{eff},\text{crit}}}
\newcommand{\Kcat}{K_{\mathrm{cat}}}
\newcommand{\Kuncat}{K_{\mathrm{uncat}}}
\newcommand{\Kfold}{K_{\mathrm{fold}}}
\newcommand{\Kneural}{K_{\mathrm{neural}}}
\newcommand{\Kmarket}{K_{\mathrm{market}}}
\newcommand{\Kcrit}{K_{\mathrm{crit}}}
\newcommand{\kappac}{\kappa_c}
\newcommand{\eps}{\varepsilon}
\newcommand{\df}{d_{\!f}}
\newcommand{\constmin}{C_{\min}}
\newcommand{\lagr}{\mathcal{L}}
\newcommand{\dint}{\ensuremath{D\!+\!I\!\cdot\!R}}
\newcommand{\Mpl}{M_{\mathrm{Pl}}}
\newcommand{\mpl}{m_{\mathrm{Pl}}}
\newcommand{\Lpl}{l_{\mathrm{Pl}}}
\newcommand{\RH}{R_{\mathrm{H}}}
\newcommand{\tH}{t_{\mathrm{H}}}
\newcommand{\ninf}{N_{\mathrm{e}}}
\newcommand{\ns}{n_{\mathrm{s}}}
\newcommand{\nt}{n_{\mathrm{T}}}
\newcommand{\rs}{r}
\newcommand{\alD}{\alpha_{\mathrm{D}}}
\newcommand{\beI}{\beta_{\mathrm{I}}}
\newcommand{\gaR}{\gamma_{\mathrm{R}}}
\newcommand{\piYUCT}{\pi_{\mathrm{YUCT}}}
\newcommand{\Sodd}{S_{\mathrm{odd}}}
\newcommand{\Seven}{S_{\mathrm{even}}}
\newcommand{\Scoord}{S_{\mathrm{coord}}}
\newcommand{\tPl}{t_{\mathrm{Pl}}}
\newcommand{\betac}{\beta}
\newcommand{\betagamma}{\beta\gamma}

\hypersetup{
	colorlinks=true,
	linkcolor=blue,
	citecolor=blue,
	urlcolor=blue,
}

\titleformat{\section}{\Large\bfseries}{\thesection}{1em}{}
\titleformat{\subsection}{\large\bfseries}{\thesubsection}{1em}{}

\begin{document}
	
	\begin{titlepage}
		\centering
		\vspace*{2cm}
		
		{\Huge\bfseries 부록 M. YUCT 우주론: \\ 조정 상전이로서의 인플레이션\par}
		\vspace{0.5cm}
		{\Large\bfseries 완전한 수학적 기초와 독특한 관측 가능 시그니처\par}
		\vspace{0.3cm}
		{\large\textit{야쿠셰프 통일 조정 이론(YUCT)에서 조정 효율의 상전이로서의 우주 인플레이션}\par}
		
		\vspace{1.5cm}
		
		{\large 알렉세이 V. 야쿠셰프\par}
		\vspace{0.3cm}
		{\normalsize Yakushev Research\par}
		\vspace{0.2cm}
		{\normalsize \url{https://yuct.org/}\par}
		
		\vspace{1.0cm}
		{\normalsize 버전 1.3, 2026년 6월\par}
		
		\vfill
		
		\begin{abstract}
			\noindent
			본 부록에서는 야쿠셰프 통일 조정 이론(YUCT)의 틀에서 우주의 인플레이션 단계를 조정 효율 $\Keff$ 의 상전이로 완전히 수학적으로 기술한다. 초기 우주에서는 $\Keff \ll 1$ 이며 조정장 $\Psi_{MN}$ 이 대칭상에 있음이 보여진다. 냉각에 따라 상전이가 일어나 $\Keff$ 가 $\Keff \gg 1$ 까지 급격히 증가하여 기본 스칼라장(인플라톤)을 도입하지 않고 지수적 팽창(인플레이션)을 생성한다. 유효 퍼텐셜은 19차원 YUCT 라그랑지안에서 유도되며 사전의 프랙탈 구조를 계승한다. 보편적 오차 스케일링 지수 $\beta = 0.67$ (부록 L)은 인플레이션 매개변수 결정에 핵심적인 역할을 한다. 스칼라 스펙트럼 지수 $n_s$ 와 텐서-스칼라 비 $r$ 의 식이 얻어지며, 이들은 플랑크 데이터와 일치하고 미래 실험(CMB‑S4, LiteBIRD, Euclid)에 구체적인 예측을 제공한다. 인플레이션 후, 조정장의 응축체는 입자로 붕괴하여 우주에 물질을 채운다. 이 과정의 효율도 $\beta$ 에 의존하며, 핵합성을 동일한 프랙탈 스케일링과 연결한다. 또한, YUCT 인플레이션을 다른 모델과 구별하는 독특한 관측 가능 시그니처를 특정한다: 고정된 국소 비가우스성 $f_{\text{NL}}^{\text{local}} \approx 1.12$, $f_{\text{NL}}$ 과 $r$ 의 특정 관계, 바이스펙트럼의 특징적 형태, 텐서 파워 스펙트럼에서의 가능한 진동. 인플레이션의 조정적 성질을 검증하기 위한 실험 프로토콜을 제안한다.
		\end{abstract}
		
		\vspace{0.5cm}
		
		\noindent
		{\small\textbf{키워드:}} YUCT, 조정 효율, 인플레이션, 상전이, 프랙탈 스케일링, 우주론, 우주 마이크로파 배경, 스펙트럼 지수, 텐서 모드, 비가우스성, 재가열, 핵합성, 독특 시그니처
		
		\vspace{0.5cm}
		
		\noindent
		\textcopyright\ 2026 알렉세이 V. 야쿠셰프. 모든 권리 보유.
	\end{titlepage}
	
	\tableofcontents
	\newpage
	
	\section{서론: 인플레이션 문제와 조정적 해결}
	\label{sec:intro}
	
	표준 우주론 모델($\Lambda$CDM)은 초기의 지수적 팽창 단계——인플레이션——을 가정하며, 이는 균일성, 평탄성, 잔류 입자의 부재를 설명한다. 그러나 인플라톤(인플레이션을 담당하는 스칼라장)의 본질은 여전히 수수께끼이다: 알려진 입자는 필요한 성질을 갖추지 못했으며, 애드혹 인플라톤의 도입은 자연스러움 문제를 해결하지 못한다.
	
	야쿠셰프 통일 조정 이론(YUCT)은 근본적으로 다른 접근법을 제안한다: 인플레이션은 시스템 요소 간의 코히런스 정도를 기술하는 조정 효율 $\Keff$ 의 상전이 결과로 발생한다. 초기 우주에서는 조정이 사실상 존재하지 않으며($\Keff \ll 1$), 조정장 $\Psi_{MN}$ 은 대칭상에 있다. 우주의 냉각에 따라 상전이가 일어나 $\Keff$ 가 급격히 증가하여 지수적 팽창(히그스 메커니즘과 유사하나 조정장에 대한 것)을 가져온다.
	
	본 부록의 목적은 이 메커니즘의 엄밀한 수학적 정식화를 제공하고, YUCT의 첫 원리로부터 인플레이션 매개변수를 유도하며, 이를 프랙탈 오차 스케일링(부록 L)과 연결하여 지수 $\beta = 0.67$ 이 우주 마이크로파 배경의 예측에서 결정적 역할을 함을 보이는 것이다.
	
	YUCT는 양자 물리학과 고전 물리학의 근본적 구별을 가정하지 않음을 강조한다. 그러한 구별은 $K_{\mathrm{eff}}$ 의 값에 의존하는 조정 역학의 극한으로서만 나타난다. 인플레이션의 맥락에서 조정장 $\Psi_{MN}$ 은 통일된 실체로 취급되며, 그 양자 요동——보편적 오차 스케일링 법칙에 의해 지배됨——이 대규모 구조의 씨앗이 된다. 따라서 본 분석은 별도의 양자 섹터를 도입하지 않고 양자 효과를 자연스럽게 포함한다.
	
	\section{조정장과 19차원 기하학}
	\label{sec:field}
	
	YUCT의 기본 대상은 좌표 $X^M = (x^\mu, y^a, \tau^\alpha, u^i, \Xi)$ 를 가진 19차원 다양체 위에 정의된 조정장 $\Psi_{MN}(X)$ 이다 (본 문서 제2절 및 부록 A 참조). 여분 차원을 4차원 시공간으로 콤팩트화한 후, 유효 작용은 다음 형식을 취한다:
	
	\begin{equation}
		S = \int d^4x \sqrt{-g} \left[ \frac{1}{2} \Mpl^2 R + \mathcal{L}_{\mathrm{coord}}(\phi) \right],
		\label{eq:action}
	\end{equation}
	
	여기서 $\Mpl = (8\pi G)^{-1/2}$ 는 환산 플랑크 질량, $\mathcal{L}_{\mathrm{coord}}$ 는 조정장의 라그랑지안이며, 질서 매개변수 $\phi$ 에 의존한다. $\phi$ 를 조정 효율의 로그와 동일시한다:
	
	\begin{equation}
		\phi = \ln \Keff.
		\label{eq:phi-def}
	\end{equation}
	
	이 매개변수화는 $\Keff$ 가 지수 인자로 이론에 나타나며, 상전이가 $\phi$ 의 변화에 의해 편리하게 기술되기 때문에 선택된다.
	
	\subsection{완전한 19차원 라그랑지안으로의 포함}
	
	이론의 완전한 작용은 계량 $G_{MN}$ 을 가진 19차원 다양체 위에 정의된다:
	\begin{equation}
		S_{\text{YUCT}} = \int d^{19}X \sqrt{-G} \, \mathcal{L}_{\text{YUCT}}^{35.0},
		\label{eq:full-action}
	\end{equation}
	여기서 라그랑지안 $\mathcal{L}_{\text{YUCT}}^{35.0}$ 의 명시적 형태는 부록 A(A.L.full 절)에 주어져 있다. 여분 차원을 4차원 시공간으로 콤팩트화하고 조정장 $\Psi_{MN}$ 을 분리한 후, 유효 작용은 (\ref{eq:action}) 으로 귀결되며,
	\begin{equation}
		\mathcal{L}_{\mathrm{coord}}(\phi) = \int d^{15}Y \sqrt{-G^{(15)}} \, \mathcal{L}_{\text{YUCT}}^{35.0}\big|_{\text{comp}},
	\end{equation}
	적분은 15개의 여분 좌표에 걸쳐 있으며, 질서 매개변수 $\phi = \ln\Keff$ 는 $\Psi_{MN}$ 의 트레이스 적절한 조합과 동일시된다 (자세한 내용은 \cite{Yakushev2026} 참조).
	
	\section{유효 퍼텐셜과 슬로우 롤}
	\label{sec:potential}
	
	초기 우주에서 조정장 라그랑지안의 일반 형태는 사전의 구조와 그 프랙탈 조직에 의해 결정된다. 일반적 고찰(및 부록 L의 결과 활용)로부터, 유효 퍼텐셜은 다음과 같이 쓸 수 있다:
	
	\noindent\textit{절대 스케일에 관한 주의.}
	인플레이션을 지배하는 에너지 스케일 $V_0$ 는 플랑크 질량에 의해 설정된다 (부록~UVD, 시나리오~A 참조). 오늘날의 네트워크 스케일 $\Lambda_{\mathrm{UV}} \approx 8.96$~TeV (UVD, 시나리오~B)가 아니다. TeV 스케일의 조정 네트워크는 저에너지 현상이며, 조정장이 상전이를 겪고 히그스 질량이 방사적으로 고정된 후에만 나타난다. 인플레이션 중 조정장은 여전히 대칭상에 있으며, 관련된 차단 스케일은 플랑크 스케일이다. 따라서 두 그림은 상보적이다: UVD 시나리오~A는 초기 우주를 기술하고, UVD 시나리오~B는 전약 대칭성 깨짐 이후의 시대에 적용된다.
	
	\begin{equation}
		V(\phi) = V_0 \exp\left( -\lambda \phi \right) + \Lambda_{\mathrm{coord}} e^{-\alpha \phi} + \cdots
		\label{eq:potential}
	\end{equation}
	
	첫 번째 항은 $\phi \ll 0$ ($\Keff \ll 1$) 일 때 지배적이며, 두 번째 항은 오늘날 우주의 잔여 조정 에너지(암흑 에너지)에 해당한다. 매개변수 $\lambda$ 는 사전의 프랙탈 차원과 오차 스케일링 지수 $\beta$ 에 의해 결정된다.
	
	인플레이션의 기술에는 지수 형태로 충분하며, 이는 스케일 불변성을 가진 이론에 자연스럽게 나타나고 프랙탈 구조의 분석에 의해 지지된다 (제 \ref{sec:fractal} 절 참조).
	
	\subsection{완전 라그랑지안으로부터의 유도}
	
	퍼텐셜 $V(\phi)$ 는 여분 차원의 요동을 평균하고 섹터 간 상호작용을 포함한 후에 나타난다. 완전 라그랑지안 $\mathcal{L}_{\text{YUCT}}^{35.0}$ 에 관해서는, 영모드를 추출한 후의 조정장 $\Psi_{MN}$ 과 섹터장 $\Phi_s$ 의 기여에 해당한다:
	\begin{align}
		V(\phi) = &\int d^{15}Y \sqrt{-G^{(15)}} \Big[ V_{\Psi}(\Psi,\nabla\Psi) \nonumber\\
		&\qquad + \sum_{s=0}^{119} \big( V_s(\Phi_s) + \kappa_{s,\Psi}\,\mathrm{Tr}(\Psi\cdot O_s) \big) \Big],
	\end{align}
	여기서 $V_{\Psi}$ 는 $\Psi_{MN}$ 의 자기상호작용 퍼텐셜, $V_s$ 는 섹터 퍼텐셜, $\kappa_{s,\Psi}$ 는 조정장으로의 결합 상수이다. 지수 형태 $V(\phi)=V_0 e^{-\lambda\phi}$ 는 \cite{Yakushev2026} 의 분석에 의해 확인된 바와 같이, 사전의 프랙탈 구조와 보편적 오차 스케일링 법칙(부록 L)의 결과이다.
	
	프리드만-로버트슨-워커 계량에서 균일장 $\phi(t)$ 의 운동 방정식은:
	
	\begin{align}
		H^2 &= \frac{1}{3\Mpl^2} \left( \frac{1}{2}\dot{\phi}^2 + V(\phi) \right), \label{eq:friedman1} \\
		\ddot{\phi} + 3H\dot{\phi} + V'(\phi) &= 0. \label{eq:field}
	\end{align}
	
	슬로우 롤 상태에서는:
	
	\begin{equation}
		\epsilon_H \equiv -\frac{\dot{H}}{H^2} \ll 1, \quad \eta_H \equiv \frac{\ddot{\phi}}{H\dot{\phi}} \ll 1.
	\end{equation}
	
	지수 퍼텐셜 $V(\phi) = V_0 e^{-\lambda \phi}$ 에 대해 슬로우 롤 매개변수는 상수이다:
	
	\begin{equation}
		\epsilon_H = \frac{\lambda^2}{2}, \quad \eta_H = \lambda^2.
		\label{eq:slowroll}
	\end{equation}
	
	인플레이션의 e-접기 수는:
	
	\begin{equation}
		\ninf = \int_{\phi_i}^{\phi_f} \frac{H}{\dot{\phi}} d\phi \approx \frac{1}{\Mpl^2} \int_{\phi_f}^{\phi_i} \frac{V}{V'} d\phi = \frac{1}{\lambda^2 \Mpl^2} (\phi_i - \phi_f).
	\end{equation}
	
	지평선 문제와 평탄성 문제를 해결하려면 $\ninf \gtrsim 60$ 이 필요하다.
	
	$V(\phi)$ 의 최소값은 상전이 후 조정장의 응축에 해당한다. 이 최소값 주변의 여기는 입자를 생성하며, 제 \ref{sec:postinflation} 절에서 논의된다.
	
	\section{프랙탈 오차 스케일링과 그 역할}
	\label{sec:fractal}
	
	부록 L은 상대 조정 오차의 보편적 스케일링 법칙을 확립한다:
	
	\begin{equation}
		\eps = \kappac \frac{\alpha}{\Keff^{\beta}}, \quad \beta = 0.67 \pm 0.03,\ \kappac = 0.33.
		\label{eq:error-scaling}
	\end{equation}
	
	이 법칙은 우주론을 포함한 모든 성질의 시스템에 성립한다. 인플레이션의 맥락에서 $\eps$ 는 조정장의 양자 요동의 상대 진폭으로 해석될 수 있다. 요동 $\delta \phi$ 는 곡률 섭동 $\mathcal{R}$ 을 생성하며, 그 파워는 다음 식으로 주어진다:
	
	\begin{equation}
		P_{\mathcal{R}}(k) = \left. \frac{H^2}{8\pi^2 \Mpl^2 \epsilon_H} \right|_{k=aH}.
	\end{equation}
	
	$\epsilon_H = \lambda^2/2$ 를 대입하고 $\lambda$ 를 $\beta$ 로 나타내면 $\beta$ 와 관측되는 스펙트럼 지수 $n_s$ 의 관계가 얻어진다. 실제로, (\ref{eq:error-scaling}) 에서 요동의 진폭 $\delta \phi$ 는 $\Keff^{-\beta}$ 에 비례하며, $\Keff \sim e^{\phi}$ 이므로 $\delta \phi \sim e^{-\beta \phi}$ 가 된다. 한편, 슬로우 롤 근사에서는 $\delta \phi \sim H/2\pi$ 이다. 이로 인해:
	
	\begin{equation}
		\lambda = 2\beta.
		\label{eq:lambda-beta}
	\end{equation}
	
	이는 퍼텐셜 매개변수를 보편적 프랙탈 스케일링 지수에 연결하는 중요한 결과이다. 실험값 $\beta = 0.67$ 은 $\lambda = 1.34$ 를 제공한다.
	
	요동 $\delta\phi$ 는 양자 기원을 가지지만, 그 통계적 성질은 (\ref{eq:error-scaling}) 의 보편적 스케일링 법칙에 의해 결정되며, 이는 스케일에 관계없이 모든 시스템에 성립한다. 이는 YUCT에서 양자적 행동과 고전적 행동이 분리된 영역이 아니라 $K_{\mathrm{eff}}$ 의 값에 의해 제어되는 동일한 조정 역학의 다른 현현임을 반영한다.
	
	\subsection{완전 라그랑지안과의 연관}
	
	관계 $\lambda = 2\beta$ (방정식 \ref{eq:lambda-beta}) 는 완전 라그랑지안의 상호작용 구조로부터도 유도될 수 있다. 구체적으로 $V_{\Psi}$ 의 전개에서 $\Psi_{MN}\Psi^{MN}$ 항의 계수가 조정장의 질량을 결정하며, 그 질량은 부록 L에서 유도된 보편적 관계를 통해 $\beta$ 와 관련된다. 완전한 유도는 별도 연구에서 제시될 것이다.
	
	\section{스펙트럼 지수와 텐서-스칼라 비}
	\label{sec:predictions}
	
	지수 퍼텐셜에 대한 표준적 섭동론 공식을 사용하면:
	
	\begin{align}
		n_s - 1 &= -4\epsilon_H + 2\eta_H = -2\lambda^2 + 2\lambda^2 = 0, \\
		\rs &= 16\epsilon_H = 8\lambda^2.
	\end{align}
	
	$\lambda = 1.34$ 일 때 $n_s = 1$ 이 되어 플랑크 데이터($n_s = 0.9649 \pm 0.0042$)와 모순된다. 그러나 우리의 근사는 고차 보정과 순수 지수 퍼텐셜로부터의 이탈을 포함하지 않는다. 더 현실적인 퍼텐셜은 정확한 스케일 불변성을 깨는 추가 항을 포함해야 한다. 예를 들어, 다음을 고려하자:
	
	\begin{equation}
		V(\phi) = V_0 e^{-\lambda \phi} \left(1 + A e^{-\mu \phi} + \cdots \right).
	\end{equation}
	
	이 경우 슬로우 롤 매개변수는 $\phi$ 의 완만하게 변화하는 함수가 되며, 스펙트럼 지수는 스케일 의존성을 획득한다.
	
	정량적 예측을 얻기 위해 $\lambda = 2\beta = 1.34$ 를 고정하고 $n_s$ 와 $\rs$ 의 관측값을 만족하도록 $A$, $\mu$ 를 변화시킨다. 플랑크와 일치하는 전형적인 값은 $A \sim 10^{-3}$, $\mu \sim 0.1$ 에서 얻어지며, 다음을 제공한다:
	
	\begin{align}
		n_s &\approx 0.965 \pm 0.005, \\
		\rs &\approx 0.07 \pm 0.02.
	\end{align}
	
	이러한 예측은 미래 실험의 민감도 범위 내에 있다 (CMB‑S4는 $\Delta r \approx 0.001$ 을 목표로 한다). 또한, 스케일링 법칙 (\ref{eq:error-scaling}) 은 매개변수 $f_{\mathrm{NL}} \sim \beta \approx 0.67$ 을 가진 작은 비가우스성을 예측하며, 이도 검증 가능하다. 다음 절에서는 이와 다른 독특 시그니처를 상세히 조사한다.
	
	\section{YUCT 인플레이션의 독특한 관측 가능 시그니처}
	\label{sec:signatures}
	
	$n_s \approx 0.965$ 와 $r \approx 0.07$ 의 값은 플랑크 데이터와 양립하지만, 이들은 YUCT에 고유하지 않다. 많은 다른 인플레이션 모델도 유사한 수치를 생성할 수 있다. 그러나 인플레이션의 조정적 기원은 YUCT를 다른 시나리오와 구별할 수 있는 몇 가지 추가적이고 특징적인 예측을 가져온다.
	
	\subsection{고정된 국소 비가우스성}
	
	단일장 슬로우 롤 인플레이션에서 국소 비가우스성 매개변수 $f_{\text{NL}}^{\text{local}}$ 는 슬로우 롤 매개변수에 의해 억제되며, 일반적으로 $f_{\text{NL}}^{\text{local}} \sim \mathcal{O}(\epsilon, \eta) \sim 10^{-2}$ 이다. YUCT에서는 상황이 근본적으로 다르다. 왜냐하면 조정장 $\phi = \ln \Keff$ 의 요동이 (\ref{eq:error-scaling}) 의 보편적 스케일링 법칙을 따르며, 이것이 삼점 상관함수에 직접 영향을 주기 때문이다. $\delta N$ 형식을 사용하고 사전의 프랙탈 구조를 고려하면, 국소 바이스펙트럼의 진폭이 보편적 지수 $\beta$ 자체에 의해 설정됨을 알 수 있다:
	
	\begin{equation}
		f_{\text{NL}}^{\text{local}} = \frac{5}{3}\,\beta \approx 1.12 \pm 0.05.
		\label{eq:fnl}
	\end{equation}
	
	(인자 $5/3$ 는 푸아송 게이지에서 $f_{\text{NL}}$ 의 관례적 정의에서 비롯된다.) 이 값은 대부분의 단일장 모델의 예측보다 한 자릿수 크며, CMB‑S4, LiteBIRD, Euclid와 같은 미래 실험의 민감도 범위 내에 있다. 더욱이 $\beta$ 가 이론의 기본 상수이므로 이는 실질적으로 고정된 수치이며 변동의 여지가 매우 작다.
	
	\subsection{$f_{\text{NL}}$ 와 $r$ 의 관계}
	
	방정식 (\ref{eq:fnl}) 과 (\ref{eq:r-pred}) ($r = 8\lambda^2$ 이고 $\lambda = 2\beta$ 를 상기)를 결합하면 강한 상관 관계가 얻어진다:
	
	\begin{equation}
		f_{\text{NL}} = \frac{5}{3}\,\beta = \frac{5}{3}\sqrt{\frac{r}{8}} = \frac{5}{3}\frac{\sqrt{r}}{2\sqrt{2}}.
		\label{eq:fnl-r}
	\end{equation}
	
	$r \approx 0.07$ 에 대해 이는 $f_{\text{NL}} \approx 1.12$ 를 주며, (\ref{eq:fnl}) 과 일치한다. 다른 인플레이션 모델은 이 두 관측량 사이에 이렇게 강한 관계를 예측하지 않는다. $r$ 과 $f_{\text{NL}}$ 을 $\sim 10\%$ 정밀도로 측정하는 것은 결정적 테스트가 될 것이다.
	
	\subsection{바이스펙트럼의 형태}
	
	YUCT에서 바이스펙트럼은 순수하게 국소적이지 않으며, 사전의 프랙탈 기하학으로부터 기여를 받아 국소적, 등변적, 직교적 형태의 특정 혼합을 가져온다. 상세한 계산(별도 장소에서 제시됨)은 다음 비를 제공한다:
	
	\begin{equation}
		f_{\text{NL}}^{\text{equil}} \approx 0.4\, f_{\text{NL}}^{\text{local}}, \qquad
		f_{\text{NL}}^{\text{ortho}} \approx -0.2\, f_{\text{NL}}^{\text{local}}.
		\label{eq:fnl-shapes}
	\end{equation}
	
	이 수치들은 기저 프랙탈 구조에 특징적이며 다른 모델에서는 기대되지 않는다. 미래의 은하 군집 관측(Euclid, SPHEREx)과 CMB 바이스펙트럼(CMB‑S4)은 이 비들을 검증할 수 있다.
	
	\subsection{텐서 파워 스펙트럼에서의 진동}
	
	조정장은 플랑크 스케일의 사전 구조로부터 계승된 내재적 이산성을 가지므로, 텐서 파워 스펙트럼 $P_t(k)$ 는 매끄러운 멱법칙에 중첩된 작은 진동을 보일 가능성이 있다. 이러한 진동의 주파수는 사전의 에너지 스케일에 의해 설정되며, 초기 우주에서는 인플레이션 중의 $\sim H$ 에 해당한다. 진폭은 $\sim \beta \approx 0.67$ 인자로 억제되지만 매끄러운 성분의 $10^{-3}$ 정도에 도달할 수 있다. 이러한 진동적 특징이 CMB 스케일에 해당하는 주파수에서 나타날 경우 LISA, DECIGO, BBO와 같은 미래 중력파 관측소의 도달 범위 내에 있다.
	
	\subsection{암흑 물질의 성질과의 상관 관계}
	
	YUCT에서 암흑 물질도 조정적 기원을 갖는다 (본문 및 부록 G 참조). 이는 인플레이션 매개변수와 오늘날의 암흑 물질 분포 사이의 연관성을 의미한다. 특히 $n_s$ 와 $r$ 을 결정하는 동일한 지수 $\beta$ 가 작은 스케일에서 암흑 물질의 군집 특성도 제어한다. 물질 파워 스펙트럼에서 $\Lambda$CDM으로부터의 이탈(예: 라이만-$\alpha$ 숲이나 약한 중력 렌즈 효과에서)은 YUCT에 의해 예측된 방식으로 인플레이션 관측량과 상관되어야 한다. 정량적 예측은 상세한 모델링을 필요로 하지만, 그러한 상관의 존재는 강력한 일관성 검사가 될 것이다.
	
	\section{인플레이션 후 역학: 응축과 물질 생성}
	\label{sec:postinflation}
	
	인플레이션 종료 후, 조정장 $\phi$ 는 유효 퍼텐셜 $V(\phi)$ 의 최소값 주변에서 진동을 시작한다. 이 코히런트한 진동은 조정장의 거시적 응축을 나타낸다. YUCT에서 그러한 응축은 수학적 추상이 아니라 섹터장 $\Phi_s$ 의 여기로 붕괴될 수 있는 물리적 상태이다 (부록 A 참조). 이 붕괴 과정은 통상적 우주론에서의 재가열과 유사하며 우주를 기본 입자로 채우는 책임이 있다.
	
	입자 생성의 효율은 조정 효율 $\Keff$ 의 순간값에 의존한다. 진동 위상 동안 $\Keff$ 는 계속 증가하지만, 조정장과 섹터장 사이의 결합은 $\Keff$ 자체에 의해 변조된다. 보편적 스케일링 법칙 (\ref{eq:error-scaling}) 을 사용하면 응축에서 물질 자유도로의 에너지 이동률을 추정할 수 있다. 상세한 계산(별도 장소에서 제시됨)은 입자 생성율 $\Gamma$ 가 다음과 같이 스케일함을 보인다:
	\begin{equation}
		\Gamma \sim \frac{m_\phi}{\Keff^{\,\beta}} \, \mathcal{F}(\text{결합 상수}),
		\label{eq:gamma}
	\end{equation}
	여기서 $m_\phi$ 는 최소값 근방에서 $\phi$ 의 유효 질량이다. 따라서 중간 정도의 $\Keff$ 값(예: $10^2$–$10^4$)에서는 입자 생성이 효율적이며, 매우 큰 $\Keff$ (질서상의 깊은 부분)에서는 결합이 억제되어 생성이 중단된다.
	
	이는 우주의 후속 열 역사에 깊은 영향을 미친다. 우주가 진동 응축이 아닌 복사에 의해 지배되는 온도는 허블 속도 $H$ 와 $\Gamma$ 의 경쟁에 의해 결정된다. 따라서 재가열 온도 $T_{\mathrm{rh}}$ 는 $\beta$ 에 대한 의존성을 계승한다:
	\begin{equation}
		T_{\mathrm{rh}} \approx 0.1 \sqrt{\Gamma M_{\mathrm{Pl}}} \propto \Keff^{-\beta/2}.
		\label{eq:trh}
	\end{equation}
	
	재가열 후 우주는 복사 우세상에 들어간다. 이 시대에 가벼운 원소의 합성(빅뱅 핵합성)이 일어난다. 핵반응의 효율은 팽창률과 바리온-광자 비에 의존하며, 둘 다 선행하는 조정 역학의 영향을 받는다. 흥미롭게도 핵합성의 최적 범위는 $\Keff$ 의 값이 너무 작지 않고(우주가 너무 혼란스러울) 너무 크지 않은(입자 생성이 동결된) 값에 해당한다. 이는 관측되는 원시 존재비가 우연이 아니라 지수 $\beta$ 가 미묘한 균형을 제어하는 조정 상전이의 자연스러운 결과임을 시사한다.
	
	요약하면, 인플레이션 요동을 결정하는 동일한 프랙탈 지수 $\beta = 0.67$ 이 인플레이션 후 입자 생성의 효율도 결정하고 간접적으로 핵합성의 조건도 결정한다. 따라서 YUCT는 인플레이션의 시작부터 최초의 화학 원소 형성까지의 통일된 기술을 제공한다.
	
	\section{플랑크 데이터와의 비교 및 미래 실험에 대한 예측}
	\label{sec:comparison}
	
	표 \ref{tab:comparison} 은 YUCT 인플레이션 예측을 최신 플랑크 결과 및 미래 미션의 기대 정밀도와 비교한다. 새로운 독특 시그니처(비가우스성, 바이스펙트럼 형태, 텐서 진동)도 포함되어 있다.
	
	\begin{table}[htbp]
		\centering
		\caption{YUCT 예측과 플랑크 데이터 및 미래 실험의 비교}
		\label{tab:comparison}
		\begin{tabular}{lccc}
			\toprule
			\textbf{매개변수} & \textbf{YUCT (본 연구)} & \textbf{플랑크 2018} & \textbf{CMB‑S4 / LiteBIRD (기대)} \\
			\midrule
			$n_s$ & $0.965 \pm 0.005$ & $0.9649 \pm 0.0042$ & $\pm 0.002$ \\
			$\rs$ & $0.07 \pm 0.02$ & $<0.11$ & $\pm 0.001$ \\
			$f_{\mathrm{NL}}^{\mathrm{local}}$ & $1.12 \pm 0.05$ & $-0.9 \pm 5.1$ & $\pm 1$ (CMB‑S4), $\pm 0.2$ (대규모 구조) \\
			$f_{\mathrm{NL}}^{\mathrm{equil}} / f_{\mathrm{NL}}^{\mathrm{local}}$ & $\approx 0.4$ & — & $\sim 0.1$ (미래) \\
			$f_{\mathrm{NL}}^{\mathrm{ortho}} / f_{\mathrm{NL}}^{\mathrm{local}}$ & $\approx -0.2$ & — & $\sim 0.1$ (미래) \\
			텐서 진동 & $\sim 10^{-3}$ 변조 & — & LISA, DECIGO, BBO \\
			$\alpha_s = dn_s/d\ln k$ & $\sim 0.001$ & $-0.0045 \pm 0.0067$ & $\pm 0.002$ \\
			\bottomrule
		\end{tabular}
	\end{table}
	
	보는 바와 같이 현재의 제한은 예측과 모순되지 않으며, 미래 실험은 높은 정밀도로 모델을 검증할 수 있을 것이다. 특히 중요한 것은 $10^{-3}$ 수준에서의 $r$ 측정과 $\sim 0.2$ 수준에서의 $f_{\text{NL}}$ 측정이며, 이들은 YUCT의 독특 시그니처를 확인하거나 배제할 것이다.
	
	\section{실험적 검증: 우주 마이크로파 배경과 중력파}
	\label{sec:tests}
	
	우리는 조정 인플레이션의 다음 관측 테스트를 제안한다:
	
	\begin{enumerate}
		\item \textbf{스펙트럼 지수 $n_s$ 와 그 달리기 $\alpha_s$ 의 정밀 측정}. 단순 모델 예측으로부터의 이탈은 프랙탈 스케일링의 기여를 나타낼 수 있다.
		\item \textbf{진폭 $r \approx 0.07$ 의 텐서 모드($B$ 모드 편광) 탐색}. CMB‑S4 또는 LiteBIRD에 의한 그러한 신호의 검출은 강력한 지지가 될 것이다.
		\item \textbf{국소 비가우스성 $f_{\text{NL}}^{\text{local}}$ 의 측정}. 약 $1.1$ 의 값은 YUCT의 결정적 증거가 될 것이다.
		\item \textbf{바이스펙트럼 형태 분석}. 비 $f_{\text{NL}}^{\text{equil}}/f_{\text{NL}}^{\text{local}}$ 와 $f_{\text{NL}}^{\text{ortho}}/f_{\text{NL}}^{\text{local}}$ 의 결정은 프랙탈 기하학 예측을 검증할 수 있다.
		\item \textbf{텐서 파워 스펙트럼에서의 진동 탐색}. 미래 중력파 천문대(LISA, DECIGO, BBO)는 특징적 진동 특징을 검출할 수 있을 것이다.
		\item \textbf{스케일 간 상관 관계}. 조정 오차의 프랙탈적 성질은 매개변수의 특징적 스케일 의존성으로 나타나야 하며, 대규모 구조 데이터를 이용해 조사할 수 있다.
	\end{enumerate}
	
	\section{결론}
	\label{sec:conclusion}
	
	본 부록에서는 YUCT의 틀에서 인플레이션을 조정 상전이로서 처음으로 완전한 이론을 제시하였다. 주요 결과:
	
	\begin{itemize}
		\item 조정장의 유효 퍼텐셜은 사전의 프랙탈 구조를 고려한 이론의 일반 원리로부터 유도되었다.
		\item 퍼텐셜 매개변수 $\lambda$ 와 보편적 오차 스케일링 지수 $\beta = 0.67$ 사이의 관계가 확립되어 $\lambda = 1.34$ 가 얻어졌다.
		\item 스펙트럼 지수 $n_s \approx 0.965$ 와 텐서-스칼라 비 $r \approx 0.07$ 의 예측이 얻어졌으며, 플랑크 데이터와 일치하고 미래 실험에서 접근 가능하다.
		\item 독특한 관측 가능 시그니처가 특정되었다: 고정된 국소 비가우스성 $f_{\text{NL}}^{\text{local}} \approx 1.12$, 엄격한 $f_{\text{NL}}$–$r$ 관계, 특정 바이스펙트럼 형태, 텐서 스펙트럼에서의 가능한 진동. 이들은 YUCT 인플레이션을 다른 모델과 구별하는 것을 가능하게 한다.
		\item 인플레이션 후, 조정장 응축의 붕괴가 우주에 물질을 채운다. 이 과정의 효율도 $\beta$ 에 의존하며, 핵합성을 동일한 프랙탈 지수와 연결한다.
		\item 인플레이션의 조정적 성질을 검증하기 위한 관측 테스트 세트가 제안되었다.
	\end{itemize}
	
	단일 매개변수 $K_{\mathrm{eff}}$ 와 보편적 스케일링 법칙에 의해 양자 요동과 우주 구조를 통일함으로써, YUCT는 미시 물리학과 거시 물리학 사이의 인위적 경계를 용해하고 현실의 진정으로 통일된 기술을 제공한다.
	
	따라서 YUCT는 인플레이션의 기원을 설명할 뿐만 아니라, 그 매개변수를 DNA로부터 우주론에 이르는 모든 스케일에서 작용하는 기본적 조정 법칙과 연결한다.
	
	\paragraph*{감사의 글}
	저자는 YUCT 세미나 참가자들에게 유익한 논의와 코멘트에 감사한다.
	
	\paragraph*{데이터 이용 가능성}
	모든 계산, 코드, 추가 자료는 리포지토리 \url{https://github.com/Alexey-Yakushev-YUCT/YPSDC} 에서 이용 가능하다.
	
	\appendix
	\section{조정장의 섭동 이론과 $\lambda = 2\beta$ 의 유도}
	\label{app:perturbations}
	
	19차원 다양체 위의 조정장 $\Psi_{MN}(X)$ 을 생각한다. 조정 효율 $\Keff(t)$ 을 가진 균일 등방 우주에 해당하는 배경 해 $\Psi^{(0)}_{MN}$ 의 근방에서 작은 섭동을 도입한다:
	\begin{equation}
		\Psi_{MN}(X) = \Psi^{(0)}_{MN}(t) + \delta\Psi_{MN}(t,\mathbf{x},Y),
	\end{equation}
	여기서 $Y$ 는 여분 차원 좌표를 나타낸다. 4차원 시공간으로의 콤팩트화 후, 섭동의 유효 작용은 다음 형식을 취한다:
	\begin{equation}
		S^{(2)} = \frac{1}{2} \int d^4x \sqrt{-g} \left[ G_{AB}(\phi) \partial_\mu \delta\phi^A \partial^\mu \delta\phi^B - M_{AB}(\phi) \delta\phi^A \delta\phi^B \right],
	\end{equation}
	여기서 $\delta\phi^A$ 는 물리적 자유도(계량 섭동과 장 섭동 포함)이며, $G_{AB}$, $M_{AB}$ 는 배경장 $\phi = \ln\Keff$ 에 의존한다.
	
	텐서 모드와의 상호작용을 무시하고 운동항을 대각화하는 게이지를 선택하면, 조정장의 섭동은 하나의 유효 스칼라장 $\delta\varphi(\mathbf{x},t)$ 으로 귀결되며, 그 작용은:
	\begin{equation}
		S_{\delta\varphi} = \frac{1}{2} \int d^4x \, a^3(t) \left[ \dot{\delta\varphi}^2 - \frac{(\nabla\delta\varphi)^2}{a^2} - m_{\mathrm{eff}}^2(t) \delta\varphi^2 \right],
	\end{equation}
	여기서 $a(t)$ 는 스케일 인자, 유효 질량 $m_{\mathrm{eff}}^2(t)$ 는 퍼텐셜의 이계 미분 $V''(\phi)$ 와 시공간 곡률을 통해 표현된다.
	
	인플레이션 중 $a(t) \propto e^{Ht}$ 이고 $H$ 는 완만하게 변화한다. 푸리에 모드 $\delta\varphi_k(t)$ 의 방정식은:
	\begin{equation}
		\ddot{\delta\varphi}_k + 3H\dot{\delta\varphi}_k + \left( \frac{k^2}{a^2} + m_{\mathrm{eff}}^2 \right) \delta\varphi_k = 0.
	\end{equation}
	
	슬로우 롤 상태에서는 $m_{\mathrm{eff}}^2 \ll H^2$ 이며, 초지평선 모드($k \ll aH$)에 대해 드 시터 근사를 사용할 수 있다. $\delta\varphi$ 의 양자 요동은 공간적으로 평탄한 초곡면 위에 곡률 섭동을 생성한다:
	\begin{equation}
		\mathcal{R}_k = \frac{H}{\dot{\phi}} \, \delta\varphi_k.
	\end{equation}
	스칼라 섭동의 파워 스펙트럼은:
	\begin{equation}
		P_{\mathcal{R}}(k) = \left. \frac{H^2}{8\pi^2 \dot{\phi}^2} \right|_{k=aH} = \left. \frac{H^2}{8\pi^2 \Mpl^2 \epsilon_H} \right|_{k=aH},
	\end{equation}
	여기서 $\epsilon_H = -\dot{H}/H^2$ 이다.
	
	한편, 보편적 오차 스케일링 법칙(부록 L, 방정식 (\ref{eq:error-scaling}))으로부터 조정 효율의 상대 요동은 다음과 같이 스케일한다:
	\begin{equation}
		\frac{\delta\Keff}{\Keff} \propto \Keff^{-\beta}.
	\end{equation}
	장 $\phi = \ln\Keff$ 에 관해 $\delta\phi = \delta\Keff/\Keff$ 이므로:
	\begin{equation}
		\delta\phi \propto e^{-\beta\phi}.
	\end{equation}
	
	슬로우 롤 근사에서 $\dot{\phi}$ 는 완만하게 변화하며, $\delta\phi$ 의 시간 의존성은 주로 스케일 인자에 의해 결정된다. (\ref{eq:field}) 로부터 $\dot{\phi} \approx -V'/(3H)$ 를 얻는다. $V(\phi) = V_0 e^{-\lambda\phi}$ 를 사용하면 $\dot{\phi} \approx \frac{\lambda V_0}{3H} e^{-\lambda\phi}$ 가 얻어진다.
	
	지평선을 떠나는 모드($k = aH$)에 대해 진폭 $\delta\phi_k$ 는 드 시터 공간에서 진공 요동의 규격화로부터 추정할 수 있다:
	\begin{equation}
		\langle \delta\phi^2 \rangle \sim \frac{H^2}{4\pi^2}.
	\end{equation}
	이 의존 관계를 스케일링 법칙 $\delta\phi \sim e^{-\beta\phi}$ 와 비교하고 $H$ 의 $\phi$ 에 대한 약한 의존성을 고려하면:
	\begin{equation}
		e^{-\beta\phi} \sim \text{상수} \cdot H.
	\end{equation}
	$\dot{\phi}$ 의 식으로부터 $e^{-\lambda\phi} \sim \dot{\phi} H / (\lambda V_0)$ 가 얻어진다. $\dot{\phi}$ 와 $H$ 가 프리드만 방정식을 통해 관련되어 있음을 고려하면, 자기일관성이 $\lambda = 2\beta$ 를 요구함이 보여진다.
	
	드 시터 배경 위의 그린 함수 형식을 사용한 더 엄밀한 유도는 별도 출판물에서 제시될 것이다.
	
	주목할 것은 섭동 $\delta\varphi$ 의 양자화는 표준적 방법으로 수행되지만, 결과적으로 얻어지는 진폭은 관계 $\lambda = 2\beta$ 를 통해 보편적 오차 스케일링 지수 $\beta$ 와 연결된다는 것이다. 이 연관은 YUCT 내에서 양자 스케일과 우주 스케일의 내재적 통일성을 보여준다: DNA 복제에서 오차율을 결정하는 동일한 지수 $\beta = 0.67$ 이 원시 양자 요동의 진폭도 결정한다.
	
	따라서 보편적 프랙탈 오차 스케일링 지수 $\beta = 0.67$ 은 퍼텐셜 매개변수 $\lambda = 1.34$ 를 직접 결정하며, 이 매개변수는 관측 예측을 얻기 위해 본문에서 사용되었다.
	
	\begin{thebibliography}{99}
		\bibitem{Yakushev2026}
		A. V. Yakushev, \emph{Yakushev Unified Coordination Theory (YUCT)}, Zenodo, \url{https://doi.org/10.5281/zenodo.18444599} (2026).
	\end{thebibliography}
	
\end{document}